% =============================================================================
% CAPÍTULO 2: PROCESSOS DE ACORDO
% =============================================================================

\aberturacapitulo{Capítulo 2}{Processos de Acordo}



\subsecao{2.1 Descrição do processo de aquisição}

O processo de aquisição (ABNT NBR ISO/IEC/IEEE 12207:2021 -- Cláusula 6.1.1) estabelece a governança da \textbf{CineSoft Solutions} para decidir entre desenvolver internamente (\emph{Make}) ou contratar/adquirir de terceiros (\emph{Buy}) os componentes, bibliotecas, APIs e serviços de infraestrutura que sustentam o produto \textbf{CineDoom} ao longo de suas três Macroentregas.


\subsubsecao{Fronteira e Justificativa de Aquisição}

A estratégia de engenharia da CineSoft Solutions concentra o esforço de desenvolvimento interno estritamente nos ativos que constituem o \textbf{núcleo de valor e diferencial competitivo} do CineDoom: o motor transacional de bilheteria e reservas com bloqueio pessimista ACID (\emph{SELECT FOR UPDATE}), a máquina de estados de assentos com expiração \emph{Just-in-Time}, o isolamento lógico multi-tenant e a camada de conformidade e anonimização da LGPD.

Em contrapartida, serviços de infraestrutura em nuvem, bancos de dados gerenciados, entrega de e-mails transacionais, retenção externa de logs, processamento financeiro PCI-DSS e mensageria fiscal SEFAZ são adquiridos de provedores especializados para reduzir o \emph{time-to-market}, otimizar o Custo Total de Propriedade (TCO) e garantir Acordos de Nível de Serviço (SLA ≥ 99,9\%).


\subsubsecao{\texorpdfstring{Análise de Aquisição (\emph{Make or Buy})}{Análise de Aquisição (Make or Buy)}}

A governança de aquisição da CineSoft Solutions consolida a decisão formal para cada elemento arquitetural do sistema, classificando a abordagem em \textbf{Sim} (solução madura/commodity adquirida pronta), \textbf{Não} (regra de negócio central desenvolvida internamente) ou \textbf{Parcial} (engine ou biblioteca externa combinada com regras de domínio proprietárias):

\tabelaabnt{Análise de Aquisição: Componentes de Software e Domínio (\emph{Make or Buy})}{CineSoft Solutions (2026)}{tab:make_or_buy_core}{{\small
\begin{tabularx}{\textwidth}{>{\fontSemiBold\color{BrandNavy}\arraybackslash}m{3.2cm} C{1.8cm} Y >{\arraybackslash}m{3.4cm}}
\arrayrulecolor{GithubTableBorder}\toprule
\rowcolor{BrandTableHeader}
\textbf{Elemento do Sistema} & \textbf{Adquirir?} & \textbf{Justificativa Técnica e Estratégica} & \textbf{Provedor / Abordagem} \\
\midrule
\arrayrulecolor{BrandGrayRule}
Motor de Bilheteria e Reserva Atômica ACID & \textbf{Não} & Regra de negócio central e diferencial competitivo; exige controle estrito sobre locks de linha (\texttt{SELECT FOR UPDATE}), ordenação anti-deadlock e expiração \emph{Just-in-Time}. & Desenvolvimento interno (CineSoft Solutions) \\
\midrule
Isolamento Multi-Tenant e Governança de Dados & \textbf{Não} & Núcleo arquitetural do modelo SaaS B2B; requer injeção determinística de \texttt{tenantId} em todas as queries e proteção contra vazamento \emph{cross-tenant}. & Desenvolvimento interno sobre Prisma Client Extensions \\
\midrule
Módulo de Criptografia e Anonimização LGPD & \textbf{Não} & Requisito legal crítico (Art.~18 e Art.~46 da LGPD) acoplado ao domínio fiscal de ingressos; combina \texttt{HMAC-SHA256} e \texttt{AES-256-GCM} nativos. & Desenvolvimento interno (\texttt{Node.js crypto}) \\
\midrule
Autenticação, Sessão JWT e Controle RBAC & \textbf{Parcial} & A gestão de tokens e cookies HTTP-only utiliza biblioteca consolidada, mas a política RBAC (\texttt{GERENTE}, \texttt{ATENDENTE}, \texttt{CLIENTE}), escopo por tenant e revogação são próprios. & \texttt{NextAuth.js} + Middleware proprietário CineSoft \\
\midrule
Interface Gráfica e Mapa Interativo de Salas & \textbf{Parcial} & Primitivas de estilo e ícones são acelerados por mercado, porém a matriz 2D reativa de poltronas e o fluxo de guichê são proprietários. & \texttt{Tailwind CSS}, \texttt{Lucide Icons} + Componentes CineSoft \\
\midrule
ORM e Validação de Schemas de Entrada & \textbf{Sim} & Bibliotecas de padrão industrial com alta segurança de tipos e prevenção nativa de \emph{SQL Injection}, cujo desenvolvimento próprio seria antieconômico. & \texttt{Prisma ORM} e \texttt{Zod} (Open Source / MIT) \\
\arrayrulecolor{GithubTableBorder}\bottomrule
\end{tabularx}
}
}

\tabelaabnt{Análise de Aquisição: Infraestrutura Cloud e Serviços Gerenciados (\emph{Make or Buy})}{CineSoft Solutions (2026)}{tab:make_or_buy_infra}{{\footnotesize
\begin{tabularx}{\textwidth}{>{\fontSemiBold\color{BrandNavy}\hsize=0.662\hsize\sloppy\arraybackslash}X C{1.69cm} >{\hsize=1.754\hsize\sloppy\arraybackslash}X >{\hsize=0.584\hsize\sloppy\arraybackslash}X}
\arrayrulecolor{GithubTableBorder}\toprule
\rowcolor{BrandTableHeader}
\textbf{Elemento do sistema} & \textbf{Adquirir pronto?} & \textbf{Justificativa} & \textbf{Possível fornecedor} \\
\midrule
\arrayrulecolor{BrandGrayRule}
Banco de Dados Relacional Gerenciado (DBaaS) & \textbf{Sim} & Elimina custos de DBA, provendo restauração pontual automatizada (PITR), \emph{connection pooling} para serverless e alta disponibilidade (≥ 99,9\%). & \textbf{Supabase Pro (PostgreSQL)} \\
\midrule
Plataforma de Hospedagem Serverless e CDN & \textbf{Sim} & Elimina a manutenção de servidores dedicados, entregando deploys atômicos com reversão para a versão anterior, proteção DDoS e distribuição global na borda. & \textbf{Vercel Pro} (Cloud PaaS) \\
\midrule
Serviço de Disparo de E-mails Transacionais & \textbf{Sim} & Evita que vouchers caiam em caixas de spam; a API especializada assegura entrega em milissegundos com autenticação SPF, DKIM e DMARC. & \textbf{Resend} (SaaS de E-mail Transacional) \\
\midrule
Observabilidade e Retenção Legal de Logs & \textbf{Sim} & Cumpre a retenção legal de logs por 1 ano (Marco Civil e LGPD) sem sobrecarregar o armazenamento do banco relacional, com busca indexada rápida. & \textbf{Axiom Cloud} (Log Drain SaaS) \\
\midrule
Gateway de Pagamento Online (Macroentrega 2) & \textbf{Sim} & Transfere a complexidade regulatória da certificação PCI-DSS Nível 1 e antifraude, automatizando a liquidação de PIX dinâmico e cartões. & \textbf{Stripe / Pagar.me} \\
\midrule
Motor de Mensageria Fiscal NFC-e (Macroentrega 3) & \textbf{Sim} & Abstrai as constantes mudanças tributárias estaduais da SEFAZ, contingência offline e assinatura digital com certificado A1, evitando passivos fiscais. & \textbf{Focus NFe / PlugNotas} \\
\arrayrulecolor{GithubTableBorder}\bottomrule
\end{tabularx}%
}
}


\subsubsecao{As Cinco Atividades do Processo de Aquisição (ISO/IEC/IEEE 12207:2017)}

A CineSoft Solutions conduz a contratação e a governança de serviços e componentes especializados por meio das cinco atividades normativas da norma:

\begin{enumerate}
\tightlist
\item \textbf{Preparação para a Aquisição:}
  \begin{itemize}
  \tightlist
  \item \emph{Descrição:} O Arquiteto de Software e a equipe técnica especificam as necessidades do sistema, volumetria estimada de transações e requisitos de segurança (TLS 1.3, criptografia de dados e conformidade com a LGPD).
  \item \emph{Saídas:} Especificação Técnica de Aquisição e Matriz de Critérios de Seleção.
  \end{itemize}
\item \textbf{Divulgação e Seleção de Fornecedores:}
  \begin{itemize}
  \tightlist
  \item \emph{Descrição:} A equipe levanta e compara soluções consolidadas para cada categoria (ex.: Neon vs.~Supabase; Resend vs.~AWS SES; Vercel vs.~AWS Amplify), pontuando maturidade do SDK, latência na América do Sul e custo escalonado.
  \item \emph{Saídas:} Parecer Comparativo de Fornecedores e Registro de Seleção.
  \end{itemize}
\item \textbf{Estabelecimento e Manutenção do Acordo:}
  \begin{itemize}
  \tightlist
  \item \emph{Descrição:} Formalização contratual via Termos de Serviço Corporativos e Aditivo de Tratamento de Dados (\emph{DPA} para a LGPD), assegurando SLA $\ge 99,9\%$, limites de requisição (\emph{rate-limiting}) e portabilidade de dados.
  \item \emph{Saídas:} Contrato de Serviço de Nuvem (DPA) e Registro de Credenciais em Cofre Seguro.
  \end{itemize}
\item \textbf{Monitoramento do Acordo:}
  \begin{itemize}
  \tightlist
  \item \emph{Descrição:} Integrado ao Controle de Projeto (Capítulo 4), acompanha-se a disponibilidade dos serviços via telemetria (latência p95/p99 e conexões do pooler), páginas públicas de status e conciliação mensal de consumo.
  \item \emph{Saídas:} Relatório Mensal de SLA/Consumo e Registro de Incidentes do Provedor.
  \end{itemize}
\item \textbf{Aceitação do Produto ou Serviço:}
  \begin{itemize}
  \tightlist
  \item \emph{Descrição:} O serviço é validado em ambiente de homologação (\emph{sandbox}) e integrado à esteira de CI/CD, exigindo aprovação em testes de estabilidade antes da liberação em produção acompanhada de plano de contingência (\emph{rollback plan}).
  \item \emph{Saídas:} Relatório de Testes em Sandbox e Termo de Liberação em Produção (\emph{Go-Live}).
  \end{itemize}
\end{enumerate}


\subsecao{2.2 Descrição do processo de fornecimento}

O processo de fornecimento (ABNT NBR ISO/IEC/IEEE 12207:2021 -- Cláusula 6.1.2) regulamenta como a \textbf{CineSoft Solutions} comercializa, implanta, licencia, mantém e suporta a plataforma \textbf{CineDoom} junto aos cinemas clientes e parceiros de negócio.


\subsubsecao{Perspectiva de Fornecimento e Ativos Comercializáveis}

Além da oferta integral da plataforma, a arquitetura modular da CineSoft Solutions permite o fornecimento autônomo de subsistemas de alto valor agregado para diferentes perfis de adquirentes:

\tabelaabnt{Perspectiva de Fornecimento e Ativos Comercializáveis}{CineSoft Solutions (2026)}{tab:perspectiva_fornecimento}{{\small
\begin{tabularx}{\textwidth}{>{\fontSemiBold\color{BrandNavy}\arraybackslash}L{3.6cm} L{3.8cm} Y}
\arrayrulecolor{GithubTableBorder}\toprule
\rowcolor{BrandTableHeader}
\textbf{Elemento com potencial} & \textbf{Possível adquirente} & \textbf{Modelo de acordo e preço} \\
\midrule
\arrayrulecolor{BrandGrayRule}
Plataforma CineDoom (Completa) & Cinemas de rua, circuitos independentes e redes de médio porte & \textbf{Assinatura SaaS B2B:} R\$\ 790/mês base + R\$\ 120/mês por sala (+ R\$\ 1.500 setup). \\
\midrule
Engine de Trava Atômica (API) & Teatros, casas de espetáculos e plataformas de eventos & \textbf{API Transacional:} R\$\ 0,08 por assento emitido (franquia mínima de R\$\ 200/mês com SLA 99,9\%). \\
\midrule
Módulo de Governança LGPD (SDK) & Software houses de varejo e entretenimento & \textbf{Licença de SDK:} R\$\ 8.900/ano (ou R\$\ 16.000 perpétua com 10h de consultoria). \\
\midrule
Portal White-Label de Reservas & Cinemas com ERP fiscal próprio de balcão & \textbf{Módulo SaaS White-Label:} R\$\ 490/mês por exibidor (domínio e SSL inclusos; R\$\ 800 setup). \\
\arrayrulecolor{GithubTableBorder}\bottomrule
\end{tabularx}
}
}


\subsubsecao{As Cinco Atividades do Processo de Fornecimento (ISO/IEC/IEEE 12207:2017)}

A CineSoft Solutions conduz o ciclo comercial e operacional de atendimento ao cliente por meio das cinco atividades normativas:

\begin{enumerate}
\tightlist
\item \textbf{Preparação do Fornecimento:}
  \begin{itemize}
  \tightlist
  \item \emph{Descrição:} A equipe de negócios e arquitetura mapeia os perfis de exibidores adquirentes, definindo a estratégia de monetização SaaS, matriz de precificação por sala e modelos de licenciamento de APIs desacopladas.
  \item \emph{Saídas:} Matriz de Planos Comerciais e Catálogo de Ativos Licenciáveis.
  \end{itemize}
\item \textbf{Resposta à Solicitação:}
  \begin{itemize}
  \tightlist
  \item \emph{Descrição:} Ao receber solicitações de cotação ou RFPs de exibidores, a CineSoft avalia a viabilidade técnica e arquitetural da rede (quantidade de salas e integração fiscal), gerando proposta com cronograma e escopo de setup.
  \item \emph{Saídas:} Proposta Técnica de Implantação e Proposta Comercial de Fornecimento.
  \end{itemize}
\item \textbf{Estabelecimento e Manutenção do Acordo:}
  \begin{itemize}
  \tightlist
  \item \emph{Descrição:} Formalização do Contrato de Licenciamento SaaS e Acordo de Tratamento de Dados (DPA/LGPD), estipulando SLA de 99,9\%, isolamento multi-tenant (\texttt{tenant\_id}), roadmap em 3 Macroentregas e condições de rescisão com portabilidade.
  \item \emph{Saídas:} Contrato de Prestação de Serviço (SaaS) e Aditivo de Privacidade e DPA.
  \end{itemize}
\item \textbf{Execução do Acordo:}
  \begin{itemize}
  \tightlist
  \item \emph{Descrição:} Provisionamento do tenant isolado, carga inicial de salas/filmes, treinamento operacional de atendentes e evolução orientada à esteira de CI/CD com controle de versão e validação nos 770 testes automatizados.
  \item \emph{Saídas:} Ambiente de Produção Provisionado e Relatório de Validação de CI/CD.
  \end{itemize}
\item \textbf{Entrega e Suporte:}
  \begin{itemize}
  \tightlist
  \item \emph{Descrição:} Condução da homologação assistida (\emph{UAT}), assinatura do Termo de Aceite de Implantação (TAI) após comprovação de zero \emph{double-booking} e ativação do suporte técnico contínuo estruturado nos 4 níveis de SLA.
  \item \emph{Saídas:} Termo de Aceite de Implantação (TAI) e Relatório Mensal de Cumprimento de SLA.
  \end{itemize}
\end{enumerate}


\subsubsecao{Termos do Acordo, Entregáveis e Evolução por Macroentregas}

O fornecimento comercial aos cinemas adquirentes contempla os seguintes entregáveis contratuais e evolução funcional:

\begin{itemize}
\tightlist
\item
  \textbf{Provisionamento de Tenant Isolado:} Criação da organização no banco de dados de produção com isolamento lógico garantido por \texttt{tenant\_id}, configuração de domínio/acesso e credenciais iniciais de gerência.
\item
  \textbf{Carga de Dados e Parametrização Inicial:} Mapeamento das salas físicas do cinema (fileiras, colunas e categorias Clássica, IMAX ou VIP) e importação da grade inicial de filmes.
\item
  \textbf{Documentação Operacional e Capacitação:} Entrega do Manual de Operação de Guichê e sessão de treinamento remoto ou presencial para atendentes e gerentes.
\item
  \textbf{Roadmap Contratual em 3 Macroentregas:}

  \begin{itemize}
  \tightlist
  \item
    \emph{Macroentrega 1 (\texttt{Release\ 1.0} --- Implantação Imediata):} Operação completa de back-office, bilheteria física atômica e portal web de reservas com pagamento presencial no balcão.
  \item
    \emph{Macroentrega 2 (\texttt{Release\ 2.0} --- Expansão Digital):} Habilitação de checkout online (PIX e cartão) e disparo de e-mails transacionais com domínio dedicado do exibidor.
  \item
    \emph{Macroentrega 3 (\texttt{Release\ 3.0} --- Automação Fiscal e Varejo):} Integração com SEFAZ para emissão automática de NFC-e, módulo de bomboniere e painel self-service de faturamento.
  \end{itemize}
\end{itemize}


\subsubsecao{Critérios de Aceite Formal pelo Cliente Adquirente}

A entrega de cada implantação ou macroentrega é considerada formalmente aceita mediante:
\begin{enumerate}
\tightlist
\item Aprovação integral da bateria de testes automatizados (770 cenários unitários, de integração e isolamento \emph{cross-tenant}) e homologação assistida (\emph{UAT}).
\item Inexistência de defeitos abertos de severidade \textbf{Crítica} (bloqueio de venda ou corrupção de dados) ou \textbf{Alta}.
\item Comprovação prática de zero ocorrência de venda dupla (\emph{double-booking}) sob disputa simultânea entre portal web e balcão físico.
\item Assinatura digital do \textbf{Termo de Aceite de Implantação (TAI)} pelo representante legal do cinema contratante.
\end{enumerate}


\subsubsecao{Acordo de Nível de Serviço (SLA) e Suporte Pós-Entrega}

A CineSoft Solutions assegura disponibilidade mensal mínima de \textbf{99,9\%} da plataforma e provê suporte contínuo estruturado em quatro modalidades de manutenção:

\tabelaabnt{Acordo de Nível de Serviço (SLA) e Suporte Pós-Entrega}{CineSoft Solutions (2026)}{tab:sla_suporte}{%
{\small
\begin{tabularx}{\textwidth}{>{\fontSemiBold\color{BrandNavy}\arraybackslash}L{3.4cm} Y C{2.2cm} >{\arraybackslash}L{3.0cm}}
\arrayrulecolor{GithubTableBorder}\toprule
\rowcolor{BrandTableHeader}
\textbf{Severidade e Modalidade} & \textbf{Descrição do Cenário} & \textbf{1ª Resposta} & \textbf{Meta de Resolução} \\
\midrule
\arrayrulecolor{BrandGrayRule}
Crítica (P1 --- Corretiva Urgente) & Indisponibilidade total da bilheteria ou falha na emissão de ingressos em horário de sessão. & Até \textbf{15 min} & Até \textbf{2 horas} (rollback em \textless{} 5 min) \\
\midrule
Alta (P2 --- Corretiva Prioritária) & Degradação no painel gerencial, falha no envio de voucher por e-mail ou lentidão acima do SLA. & Até \textbf{1 hora} & Até \textbf{8 horas} \\
\midrule
Normal (P3 --- Adaptativa / Preventiva) & Ajustes de compatibilidade de navegador, dúvidas operacionais ou otimização preventiva de índices. & Até \textbf{4h úteis} & Até \textbf{24 horas úteis} \\
\midrule
Evolutiva (P4 --- Perfectiva) & Solicitação de novos relatórios, filtros ou melhorias de usabilidade no Portal do Cliente. & Até \textbf{2 dias úteis} & Planejada na próxima \emph{Sprint / Release} \\
\arrayrulecolor{GithubTableBorder}\bottomrule
\end{tabularx}%
}
}


\subsecao{2.3 Identificação dos fornecedores e parceiros relevantes}

A operação contínua e comercial da plataforma CineDoom apoia-se em provedores de tecnologia de primeira linha. A tabela abaixo detalha cada fornecedor crítico, seu papel arquitetural, a estrutura de contratação escalonada entre a validação inicial (Macroentrega 1) e a operação comercial plena (Macroentregas 2 e 3), bem como os impactos em disponibilidade, custo e segurança:

\tabelaabnt{Identificação dos Fornecedores e Parceiros Relevantes}{CineSoft Solutions (2026)}{tab:fornecedores_parceiros}{%
{\small
\begin{tabularx}{\textwidth}{>{\fontSemiBold\color{BrandNavy}\arraybackslash}L{2.8cm} L{3.2cm} L{3.2cm} Y}
\arrayrulecolor{GithubTableBorder}\toprule
\rowcolor{BrandTableHeader}
\textbf{Fornecedor} & \textbf{Categoria e serviço} & \textbf{Plano e custo comercial} & \textbf{Impacto e mitigação de riscos} \\
\midrule
\arrayrulecolor{BrandGrayRule}
Vercel Inc. & PaaS \& CDN Edge (Next.js 14) & \textbf{US\$\ 20/mês} por dev (Plano Pro; piloto Hobby grátis) & SLA 99,99\% na borda. Mitigação: build determinístico e rollback em 1 clique. \\
\midrule
Supabase / Neon & DBaaS PostgreSQL (ACID) & \textbf{US\$\ 25 a US\$\ 69/mês} (Plano Pro) & Crítico para locks \texttt{SELECT FOR UPDATE}. Mitigação: réplicas e migrações Prisma. \\
\midrule
Resend & E-mail transacional (API) & \textbf{US\$\ 20/mês} (50.000 e-mails/mês) & Impacta envio de vouchers (\texttt{CA-11}). Mitigação: exibição imediata em tela. \\
\midrule
Axiom Inc. & Log drain e auditoria (LGPD) & \textbf{US\$\ 25/mês} (1 TB/mês; retenção 365 dias) & Rastreabilidade legal. Mitigação: mascaramento de PII no logger. \\
\midrule
GitHub Inc. & CI/CD e esteira Actions & \textbf{US\$\ 4/usuário/mês} (Plano Team) & Sustenta testes automatizados. Mitigação: bloqueio de PR sem aprovação. \\
\midrule
Stripe / Pagar.me \emph{(M2)} & Gateway financeiro PCI-DSS & \textbf{0,99\% PIX / 3,99\% + R\$\ 0,39 Cartão} & Risco de fraude. Mitigação: webhooks idempotentes e pagamento balcão. \\
\midrule
Focus NFe \emph{(M3)} & Mensageria fiscal NFC-e & \textbf{R\$\ 149/mês} (até 2.000 notas/mês) & Conformidade tributária. Mitigação: fila assíncrona com contingência offline. \\
\arrayrulecolor{GithubTableBorder}\bottomrule
\end{tabularx}%
}
}
